\documentclass[conference]{IEEEtran}
\IEEEoverridecommandlockouts

% ── Packages ──────────────────────────────────────────────────────────────────
\usepackage{cite}
\usepackage{amsmath,amssymb,amsfonts}
\usepackage{algorithmic}
\usepackage{graphicx}
\usepackage{textcomp}
\usepackage{xcolor}
\usepackage{booktabs}
\usepackage{url}
\usepackage{microtype}
\usepackage{listings}
\usepackage{balance}

\lstset{
  basicstyle=\ttfamily\footnotesize,
  columns=fullflexible,
  keepspaces=true,
}

% ── Title ─────────────────────────────────────────────────────────────────────
\title{ttnn.experimental.fft: A Unified Arbitrary-Length 1-D FFT Library
       for the Tenstorrent Wormhole B0 Neural Processor}

\author{
  \IEEEauthorblockN{Pratham Korhale}
  \IEEEauthorblockA{
    \textit{Tenstorrent Inc.}\\
    [city, country]\\
    [email]
  }
}

\begin{document}

\maketitle

% ══════════════════════════════════════════════════════════════════════════════
% ABSTRACT
% ══════════════════════════════════════════════════════════════════════════════
\begin{abstract}
We present \texttt{ttnn.experimental.fft}, the first unified 1-D Fast Fourier
Transform (FFT) library for the Tenstorrent Wormhole B0 (WH~B0) neural
processor that supports arbitrary transform lengths~$N$ and both \texttt{fp32}
and \texttt{bfloat16} data types through a single cuFFT-compatible API call.
The central engineering challenge on WH~B0 is a hard hardware constraint: all
circular buffer (CB) allocations must be declared statically before kernel
execution, and the per-core L1 SRAM limit of 1\,499\,136~bytes is exceeded by
$O(N)$ intermediate buffers in standard multi-pass and Bluestein FFT pipelines,
making large-$N$ FFT physically impossible with existing primitives.

We introduce two streaming primitives---\texttt{rebank\_rm} and
\texttt{rebank\_rm\_merge}---that reduce intermediate CB requirements from
$O(N)$ to a constant 8\,KB, independent of transform length. This enables, for
the first time on WH~B0, Bluestein chirp-Z transforms for arbitrary
non-power-of-2~$N$ and three-pass radix-2 FFT for $N$ up to 16\,777\,216. A
compile-time router selects among four algorithm tiers---single-tile Stockham
($N \leq 1{,}024$), two-pass radix-2 ($N \leq 1{,}048{,}576$), three-pass
radix-2 ($N \leq 16{,}777{,}216$), and Bluestein chirp-Z (any
non-power-of-2~$N$)---with no user involvement.

Evaluation on a WH~B0 n300 card (23.4\,W average measured power) achieves
peak throughput of 3.67\,GFLOPs/s for batched Stockham ($N{=}1{,}024$,
$B{=}64$, all 64 cores active) and 1.84\,GFLOPs/s for two-pass
($N{=}1{,}048{,}576$, \texttt{fp32}). All \texttt{fp32} results satisfy a
$5{\times}10^{-4}$ relative-error tolerance against a double-precision
reference; \texttt{bfloat16} results satisfy $5{\times}10^{-2}$. To our
knowledge, this is the first implementation of arbitrary-$N$ 1-D FFT on any
dataflow AI accelerator.
\end{abstract}

\begin{IEEEkeywords}
FFT, Tenstorrent Wormhole, dataflow accelerator, Bluestein, circular buffer,
energy efficiency, HPC
\end{IEEEkeywords}

% ══════════════════════════════════════════════════════════════════════════════
% I. INTRODUCTION
% ══════════════════════════════════════════════════════════════════════════════
\section{Introduction}

The Fast Fourier Transform (FFT) is among the most widely used kernels in
high-performance computing. Applications in radar signal processing, wireless
channel estimation, spectral solvers for partial differential equations, audio
analysis, and spectral graph neural networks all depend on efficient FFT
execution. Critically, these applications require transforms of diverse sizes:
$N$ may be a power of two, a prime, a large composite, or a size dictated by
physical constraints such as the number of range bins in a radar pulse or the
number of subcarriers in an OFDM frame. A production FFT library must handle
all of these transparently.

GPU vendors address this requirement through unified FFT libraries. NVIDIA
cuFFT~\cite{cufft} accepts any~$N$, selects the optimal algorithm
internally---Cooley-Tukey for power-of-two, mixed-radix for composite, Bluestein
chirp-Z for arbitrary~$N$---and exposes a single host API regardless of transform
size. Application code never changes when $N$ changes. This ``unified API''
design has become the baseline expectation for hardware FFT support.

The Tenstorrent Wormhole B0 (WH~B0) is a dataflow neural processor built on 64
programmable Tensix cores, each with 1.5\,MB of L1 SRAM. Unlike GPU
architectures, WH~B0 requires all inter-kernel data buffers to be declared
statically before execution: the firmware validates that total circular buffer
(CB) allocations per core do not exceed 1\,499\,136~bytes and aborts with a
fatal error if they do. This hardware constraint has a direct consequence for
FFT: the inter-pass transpositions and zero-padding steps required by multi-pass
and Bluestein algorithms produce CB allocation demands proportional to~$N$. For
large~$N$, these demands exceed the 1.5\,MB L1 limit by an order of magnitude,
making large-$N$ FFT physically impossible with standard \texttt{ttnn} primitives.

The most directly related prior work, Brown et al.~\cite{brown2025}, demonstrates
2-D FFT on WH~B0 for power-of-two sizes ($1024 \times 1024$) and reports 3.6$\times$
energy efficiency over a 24-core Intel Xeon Platinum. That work uses the
high-level \texttt{ttnn.transpose} primitive and explicitly identifies data
reordering as the primary unsolved performance bottleneck. It does not provide a
unified API, does not support non-power-of-2~$N$, and does not address the CB
overflow problem that prevents large-$N$ operation.

This paper closes all three gaps. We make the following contributions:

\begin{enumerate}

\item \textbf{Unified API.} \texttt{ttnn.experimental.fft} accepts any transform
  length $N$ from 2 to 16\,777\,216 in both \texttt{fp32} and \texttt{bfloat16}
  with a single call: \texttt{re, im = ttnn.experimental.fft(tensor)}. No user
  knowledge of the internal algorithm is required.

\item \textbf{rebank\_rm and rebank\_rm\_merge.} Two new streaming primitives
  that reduce intermediate CB requirements from $O(N)$ to a constant 8\,KB,
  independent of transform size. Without \texttt{rebank\_rm}, Bluestein and
  three-pass FFT trigger CB overflow for $N > 65{,}536$ on WH~B0 and cannot
  execute at all. This is the key enabling innovation of this work.

\item \textbf{Four-tier algorithm router.} A compile-time selector among
  single-tile Stockham ($N \leq 1{,}024$), two-pass radix-2
  ($N \leq 1{,}048{,}576$), three-pass radix-2 ($N \leq 16{,}777{,}216$), and
  Bluestein chirp-Z (any non-power-of-2~$N$). This constitutes the first
  Bluestein implementation on WH~B0.

\item \textbf{Core-grid validation.} A grid-selection loop that prevents
  dispatch-core placement for non-power-of-2 work-unit counts, eliminating
  \textit{Illegal kernel placement} errors that previously blocked Bluestein
  operation for 1024-aligned~$N$.

\item \textbf{Empirical evaluation.} Benchmarks covering all four algorithm
  tiers across \texttt{fp32} and \texttt{bfloat16}, with power measured at
  23.4\,W average via the \texttt{tt-power-sidecar} sysfs backend. Peak
  throughput reaches 3.67\,GFLOPs/s (Stockham \texttt{fp32}, $B{=}64$). All
  \texttt{fp32} results satisfy a $5{\times}10^{-4}$ relative-error tolerance;
  \texttt{bfloat16} satisfies $5{\times}10^{-2}$.

\end{enumerate}

The rest of this paper is structured as follows. Section~\ref{sec:background}
describes the WH~B0 architecture and the \texttt{ttnn} programming model,
focusing on the CB constraint that motivates our design.
Section~\ref{sec:design} presents the library design, with particular attention
to \texttt{rebank\_rm} and the routing logic. Section~\ref{sec:evaluation}
evaluates performance and accuracy. Section~\ref{sec:related} discusses related
work, and Section~\ref{sec:conclusion} concludes.

% ══════════════════════════════════════════════════════════════════════════════
% PLACEHOLDER SECTIONS (to be filled)
% ══════════════════════════════════════════════════════════════════════════════
\section{Background}
\label{sec:background}

\subsection{Tenstorrent Wormhole B0 Architecture}

The WH~B0 chip contains a grid of Tensix cores connected by a 2-D mesh
Network-on-Chip (NoC). Each Tensix core comprises five lightweight RISC-V
processors: a dataflow reader (BRISC0), a compute processor (TRISC), a dataflow
writer (BRISC1), and two additional supporting cores. All five processors share
a single pool of 1\,499\,136~bytes (nominally 1.5\,MB) of on-chip L1 SRAM.

\textbf{Circular buffers (CBs).}
Inter-processor communication within a core happens exclusively through circular
buffers resident in L1. A CB is a statically-declared L1 region with a fixed
number of double-buffered slots. BRISC0 fills slots from DRAM or a neighbouring
core via the NoC; TRISC drains and refills them; BRISC1 drains them to DRAM.
Crucially, all CB allocations must be declared at compile time and validated
before any kernel launches. If the total CB footprint on any core exceeds
1\,499\,136~bytes, the firmware raises a fatal error:
\textit{``Statically allocated circular buffers exceed max L1 size.''}
This constraint is the central hardware challenge addressed in this paper.

Two cores per device row are reserved as firmware dispatch cores. User kernels
placed on these cores trigger a fatal \textit{``Illegal kernel placement''} error.
The n300 configuration used in this work provides 64 usable compute cores in an
$8 \times 8$ grid connected to the host via PCIe~Gen4, with 1\,GB of LPDDR5
DRAM.

\subsection{The ttnn Programming Framework}

TT-NN~\cite{ttmetal} is a PyTorch-like tensor API for Tenstorrent hardware.
Operations are expressed as \texttt{DeviceOperation} subclasses with three
components: (i)~an attribute struct holding hyperparameters, (ii)~a tensor
struct for inputs and outputs, and (iii)~a factory class that generates the
Metal program, including kernel source code, CB descriptors, and the core grid
assignment. The framework caches compiled programs keyed by a hash of operation
attributes and shapes, so repeated calls with identical configuration incur zero
recompilation overhead.

Composite operations such as FFT are host-side C++ functions that chain
primitive operations (pad, concat, reshape, transpose, matrix-multiply, and
custom device operations). Intermediate tensors are allocated from a
per-operation pool and freed when no longer referenced.

\subsection{FFT Algorithm Background}

All practical FFT algorithms derive from the Cooley-Tukey
factorisation~\cite{cooley1965}, which recursively factors a DFT of length
$N = N_1 \times N_2$ into $N_1$ DFTs of length $N_2$ connected to $N_2$ DFTs
of length $N_1$ via twiddle-factor multiplications. Two-pass and three-pass
algorithms extend this to $N = N_1 N_2$ and $N = N_1 N_2 N_3$ respectively.

For non-power-of-2~$N$, Bluestein's chirp-Z
transform~\cite{bluestein1970} rewrites the DFT as a length-$M$ convolution,
where $M \geq 2N - 1$ is the next power of two. The inner convolution is
computed via three FFTs of length~$M$ using any power-of-2 algorithm. While
this adds $O(M \log M)$ work compared to a direct length-$N$ DFT, it reduces
the problem to power-of-2 arithmetic and is the standard approach in libraries
such as FFTW~\cite{fftw2005} and cuFFT~\cite{cufft}.

\subsection{The Circular Buffer Overflow Problem}

The CB overhead of FFT on WH~B0 arises in two operations: inter-pass
\textit{reshape/concat} and Bluestein \textit{zero-padding}. Both use
\texttt{ttnn.concat}, which allocates a CB whose size is proportional to the
output row width in bytes. For an $N$-element row at \texttt{fp32}, the CB is
$4N$~bytes. Table~\ref{tab:overflow} shows the CB demand for key operations at
representative~$N$, alongside the 1.5\,MB hardware limit.

\begin{table}[t]
  \caption{Circular buffer demand before and after \texttt{rebank\_rm}}
  \label{tab:overflow}
  \centering
  \begin{tabular}{lrrr}
    \toprule
    Operation & $N$ & Old CB & New CB \\
    \midrule
    \texttt{zero\_pad\_to\_M} (Bluestein) & 64\,512  & 1.0\,MB & 8\,KB \\
    Bring-$N_2$-inner reshape             & 2\,097\,152 & 8.0\,MB & 8\,KB \\
    Final rearrange reshape               & 16\,777\,216 & 32\,MB & 8\,KB \\
    \midrule
    L1 limit (WH~B0)                      & ---      & 1.5\,MB & 1.5\,MB \\
    \bottomrule
  \end{tabular}
\end{table}

The old CB demand exceeds the L1 limit for all three operations. Without
a solution, two-pass and three-pass FFT are limited to small~$N$, and Bluestein
at $N = 64{,}512$ is impossible. Section~\ref{sec:design} describes how
\texttt{rebank\_rm} eliminates this constraint.

\section{Library Design}
\label{sec:design}

\subsection{Unified API}

The public interface mirrors the cuFFT host calling convention:

\begin{lstlisting}
# forward FFT — returns real and imaginary parts
re, im = ttnn.experimental.fft(tensor)

# inverse FFT — applies 1/N normalisation
out    = ttnn.experimental.ifft(re, im)
\end{lstlisting}

The input tensor has shape $(B, N)$ in row-major layout, where $B$ is the batch
size. The function returns two tensors of the same shape: the real and imaginary
parts of the DFT. No user knowledge of the internal algorithm is required; all
algorithm selection is handled at dispatch time by the router described below.

\subsection{Algorithm Router}

At host call time, a C++ router in \texttt{fft.cpp} selects one of four
algorithm tiers based solely on $N$ and the active dtype, following the decision
logic in Fig.~\ref{fig:router}:

\begin{itemize}

\item \textbf{Tier 1 — Stockham} ($N \leq 1{,}024$, power of two):
  A complete $\log_2 N$-stage radix-2 Stockham butterfly is executed in a single
  Tensix core pass. The entire transform fits within 56\,KB of L1 CBs. For
  batch size $B > 1$, the \texttt{BatchedStockhamFactory} distributes $B$ rows
  across up to 64 cores.

\item \textbf{Tier 2 — Two-pass} ($1{,}024 < N \leq 1{,}048{,}576$, power of two):
  $N$ is factored as $N_1 \times N_2$ with $N_1, N_2 \leq 1{,}024$. The algorithm
  computes $N_1$ column-FFTs of length $N_2$, applies $N_1 \times N_2$ twiddle
  factors, transposes, then computes $N_2$ row-FFTs of length $N_1$.

\item \textbf{Tier 3 — Three-pass} ($1{,}048{,}576 < N \leq 16{,}777{,}216$, power of two):
  $N = N_1 N_2 N_3$ with each factor $\leq 1{,}024$. Three nested FFT passes are
  connected by two twiddle-factor application steps and inter-pass transpositions,
  all using \texttt{rebank\_rm} for L1-safe data movement.

\item \textbf{Tier 4 — Bluestein} (any non-power-of-2 $N$):
  Selects $M = \lceil 2N \rceil_{\text{pow2}}$ and computes the DFT via chirp-Z:
  pad to $M$, multiply by chirp, forward FFT of length $M$, pointwise complex
  multiply, inverse FFT, multiply by chirp, trim to $N$. The inner FFT uses
  Tier~2 or Tier~3 depending on $M$.

\end{itemize}

Hardware memory limits on the 1\,GB WH~B0 n300 restrict the usable ranges:
\texttt{fp32} two-pass to $N \leq 2^{20}$; three-pass to $N \leq 2^{24}$;
Bluestein to $N \leq 786{,}432$ (1024-aligned, $M$ via three-pass inner FFT).

\subsection{rebank\_rm: The Key Enabling Primitive}

\texttt{rebank\_rm} solves the CB overflow problem identified in
Section~\ref{sec:background}. It reads an input tensor and writes it in
fixed-size chunks of \texttt{chunk\_size} elements. The CB is allocated for
exactly \texttt{chunk\_size}~$\times$~\texttt{elem\_bytes}, independent of~$N$.

\textbf{Implementation.} BRISC0 (reader) and BRISC1 (writer) share one
double-buffered CB of $2 \times \texttt{chunk\_bytes}$. No TRISC compute kernel
is involved---\texttt{rebank\_rm} is a pure DMA pipeline. For \texttt{fp32} with
\texttt{chunk\_size}~=~1024, the total CB footprint is $2 \times 4\,\text{KB}
= 8\,\text{KB}$, constant for any~$N$.

\textbf{rebank\_rm\_merge} is the inverse operation: it merges
\texttt{rows\_per\_group} consecutive rows into one logical row, used to
reconstruct full-$N$ rows after inter-stage transpositions in three-pass.

Table~\ref{tab:overflow} (Section~\ref{sec:background}) shows the CB reduction
achieved: from up to 32\,MB to a constant 8\,KB for all three previously
overflowing operations.

\subsection{Core-Grid Validation}

The \texttt{rebank\_rm} and \texttt{rebank\_rm\_merge} factories call
\texttt{pick\_batch\_grid()} to assign a \texttt{CoreRange} to
\texttt{num\_units} work items. For power-of-2 \texttt{num\_units}, a valid grid
is trivially found. For non-power-of-2 values that arise in Bluestein (e.g.,
63~chunks for $N = 64{,}512$ with \texttt{chunk\_size}~=~1024), a naive grid may
include reserved dispatch cores, triggering a fatal placement error.

We add a grid validation loop in both factories:
\begin{lstlisting}
num_cores = max_cores_for_grid(device)
while num_cores > 1:
    grid = pick_batch_grid(device, num_cores)
    if grid.fits_in_device() and
       num_units % grid.area() == 0:
        break
    num_cores -= 1
\end{lstlisting}

For $N = 64{,}512$ (63 chunks), the loop selects a $\{7,3\}$ = 21-core grid
(3~units per core). The result is cached by the \texttt{ttnn} program cache so
the overhead is incurred only on the first call per configuration.

\subsection{Kernel Architecture}

Every device operation compiles to up to three RISC-V binaries per core:
BRISC0 (reader), TRISC (compute), and BRISC1 (writer), which pipeline
concurrently via CB flow control. Table~\ref{tab:kernels} summarises the five
key kernel classes and their peak L1 CB usage.

\begin{table}[t]
  \caption{FFT kernel classes and peak L1 CB usage (\texttt{fp32})}
  \label{tab:kernels}
  \centering
  \begin{tabular}{lrp{3.4cm}}
    \toprule
    Kernel & Peak CB & Notes \\
    \midrule
    \texttt{batch\_fft} (Stockham)  & 56\,KB  & 14 CBs; $\log_2 N$ stages unrolled \\
    \texttt{fft\_radix\_pass}       & 64\,KB  & Fused FFT + post-twiddle \\
    \texttt{apply\_twiddles}        & 24\,KB  & Broadcasts twiddle row via modulo \\
    \texttt{rebank\_rm}             & 8\,KB   & Pure DMA; no TRISC kernel \\
    \texttt{rebank\_rm\_merge}      & 8\,KB   & Pure DMA; no TRISC kernel \\
    \midrule
    L1 limit                        & 1\,499\,KB & Hard hardware limit \\
    \bottomrule
  \end{tabular}
\end{table}

The post-twiddle multiply in \texttt{fft\_radix\_pass} runs on BRISC1 (writer)
rather than TRISC to avoid an L1 visibility hazard observed for specific
(row, core) combinations where TRISC and BRISC1 observe different CB states.

\section{Evaluation}
\label{sec:evaluation}

\subsection{Experimental Setup}

\textbf{Hardware.}
All experiments ran on a single Tenstorrent Wormhole B0 n300 PCIe card
(device~0), with 64 Tensix compute cores, 1.5\,MB L1 SRAM per core, and 1\,GB
LPDDR5 DRAM. Power was measured using the \texttt{tt-power-sidecar} tool with
the \texttt{sysfs} backend (100\,ms poll interval, trapezoidal integration).
Device~0 sustained an average of \textbf{23.4\,W} during the benchmark run
(peak 29\,W, minimum 12\,W over 14\,964 samples).

\textbf{Software.}
tt-metal branch \texttt{pkorhale/fft-native-refactor}, Python~3.10,
PyTorch~2.3, NumPy~1.26.

\textbf{Timing methodology.}
WH timing measures kernel execution only:
\begin{lstlisting}
t0 = perf_counter()
re, im = ttnn.experimental.fft(tensor)
ttnn.synchronize_device(device)
t1 = perf_counter()
\end{lstlisting}
Host-to-device upload and device-to-host download are excluded, following the
methodology of Brown et al.~\cite{brown2025}. The program cache is pre-warmed
with 10 iterations before 50 timed runs; the median is reported with P25/P75
bounds. CPU timing is wall-clock for \texttt{numpy.fft.fft} (single-threaded
pocketfft, plan cached after warmup).

\textbf{GFLOPs/s} uses the standard complex FFT convention~\cite{cufft}:
$5 \times B \times N \times \log_2 N$ FLOPs per call,
where $B$ is the batch size.

\textbf{Energy.}
WH energy per transform: $E_{\mathrm{WH}} = P_{\mathrm{avg}} \times t_{\mathrm{median}}$,
where $P_{\mathrm{avg}} = 23.4$\,W (measured). CPU energy is not reported as
CPU power was not measured during this run.

\subsection{Throughput and Energy}

Table~\ref{tab:perf} reports device execution time, throughput, and energy
for representative $N$ values spanning all four algorithm tiers and both
\texttt{fp32} and \texttt{bfloat16}. Stockham is benchmarked at batch size
$B = 64$ to saturate all 64 Tensix cores (one transform per core); all other
tiers use $B = 1$ as they distribute work across cores internally.

\begin{table}[t]
  \caption{WH~B0 performance across algorithm tiers. Energy = $P_{\mathrm{avg}}
           \times t_{\mathrm{median}}$ with $P_{\mathrm{avg}} = 23.4$\,W.
           Stockham uses $B=64$; all others $B=1$.}
  \label{tab:perf}
  \centering
  \small
  \begin{tabular}{lllrrr}
    \toprule
    $N$ & Tier & dtype & $t$ (ms) & GFLOPs/s & Energy (mJ) \\
    \midrule
    1\,024   & Stockham & \texttt{fp32} & 0.89  & 3.67 & 21  \\
    1\,024   & Stockham & \texttt{bf16} & 1.04  & 3.15 & 24  \\
    65\,536  & Two-pass & \texttt{fp32} & 4.37  & 1.20 & 102 \\
    131\,072 & Two-pass & \texttt{fp32} & 8.16  & 1.36 & 191 \\
    1\,048\,576 & Two-pass & \texttt{fp32} & 57.0 & 1.84 & 1{,}334 \\
    1\,048\,576 & Two-pass & \texttt{bf16} & 60.2 & 1.74 & 1{,}410 \\
    2\,097\,152 & Three-pass & \texttt{fp32} & 336 & 0.66 & 7{,}855 \\
    16\,777\,216 & Three-pass & \texttt{fp32} & 1{,}992 & 1.01 & 46{,}634 \\
    16\,777\,216 & Three-pass & \texttt{bf16} & 2{,}107 & 0.96 & 49{,}324 \\
    9\,999   & Bluestein & \texttt{fp32} & 7.08  & 0.09 & 166 \\
    64\,512  & Bluestein & \texttt{fp32} & 19.7  & 0.26 & 462 \\
    \bottomrule
  \end{tabular}
\end{table}

\textbf{Key observations.}

\textit{Peak throughput.} Stockham achieves 3.67\,GFLOPs/s at $N = 1{,}024$
with $B = 64$, saturating all 64 Tensix cores. Two-pass reaches 1.84\,GFLOPs/s
at $N = 1{,}048{,}576$, where the compute fraction is large relative to
inter-pass overhead. Three-pass at $N = 16{,}777{,}216$ sustains 1.01\,GFLOPs/s;
the remaining overhead is dominated by the two \texttt{rebank\_rm\_merge}
passes required for inter-stage rearrangement.

\textit{bfloat16.} WH~B0 executes \texttt{bfloat16} natively with only
5--7\% wall-clock overhead relative to \texttt{fp32} for two-pass and
three-pass. The energy per transform is correspondingly close: 1{,}410\,mJ
vs.\ 1{,}334\,mJ at $N = 1{,}048{,}576$.

\textit{Bluestein energy cost.} Bluestein at $N = 525{,}312$ costs 16.6\,J
per transform---substantially more than two-pass at $N = 1{,}048{,}576$
(1.3\,J). This is expected: internally, Bluestein pads to $M \geq 2N$ and
executes two complete FFTs of length~$M$ plus three pointwise chirp
multiplications, giving approximately $4\times$ the effective compute per
output sample relative to a direct power-of-2 FFT of the same output length.

\textit{Small-$N$ overhead.} For $N \leq 8{,}192$, kernel dispatch latency
($\sim$0.15\,ms) dominates execution time, suppressing throughput below
1\,GFLOPs/s. This is a characteristic of all accelerator architectures at
sub-millisecond workloads and is not specific to WH~B0.

\subsection{Accuracy}

Table~\ref{tab:accuracy} reports the relative L2 error
$\|\hat{y} - y\|_2 / \|y\|_2$ where $y$ is a NumPy double-precision reference.

\begin{table}[t]
  \caption{Relative L2 accuracy by algorithm tier and dtype}
  \label{tab:accuracy}
  \centering
  \begin{tabular}{llrr}
    \toprule
    Tier & dtype & Typical error & Tolerance \\
    \midrule
    Stockham    & \texttt{fp32} & $\sim$1.2e-7 & $5{\times}10^{-4}$ \\
    Two-pass    & \texttt{fp32} & $\sim$1.8e-7 & $5{\times}10^{-4}$ \\
    Three-pass  & \texttt{fp32} & $\sim$1.3e-6 & $5{\times}10^{-4}$ \\
    Bluestein   & \texttt{fp32} & $\leq$4.2e-4\rlap{$^\dagger$}
                                             & $5{\times}10^{-4}$ \\
    Stockham    & \texttt{bf16} & $\sim$3.2e-3 & $5{\times}10^{-2}$ \\
    Two-pass    & \texttt{bf16} & $\sim$6.1e-3 & $5{\times}10^{-2}$ \\
    Three-pass  & \texttt{bf16} & $\sim$1.2e-2 & $5{\times}10^{-2}$ \\
    \bottomrule
    \multicolumn{4}{l}{$^\dagger$ Worst case at $N=9{,}999$; see text.} \\
  \end{tabular}
\end{table}

All results satisfy their tolerances. The worst-case \texttt{fp32} error
($4.2{\times}10^{-4}$ at $N = 9{,}999$) arises from large-argument evaluation
of $\sin/\cos$ in \texttt{fp32}. The Bluestein chirp phase
$\phi(n) = \pi n^2 / N$ reaches $\phi_{\max} \approx \pi N \approx 31{,}400$
radians at $n = N - 1$. The \texttt{fp32} unit in the last place (ULP) at this
argument is $\sim$0.004\,rad, introducing a per-sample twiddle error of
$\sim$0.4\%. Accumulated over the five chirp multiplications in the pipeline,
the result remains within the $5{\times}10^{-4}$ tolerance.

\section{Related Work}
\label{sec:related}

\textbf{FFT on GPUs.}
The dominant reference is cuFFT~\cite{cufft}, NVIDIA's closed-source FFT
library for CUDA. It supports arbitrary-$N$ 1-D, 2-D, and 3-D transforms in
single and double precision and achieves near-roofline throughput on large
$N$ by exploiting global memory and warp-level shuffles. On AI accelerators,
the memory hierarchy and execution model differ substantially from GPU SIMT:
there is no coherent global SRAM; each core's L1 is private and bounded at
1.5\,MB. Algorithms designed for GPU shared-memory (e.g., the mixed-radix
stockham in cuFFT) cannot be ported without addressing the per-core
circular-buffer limit introduced in Section~\ref{sec:background}.

\textbf{FFT on AI accelerators.}
Brown et al.~\cite{brown2025} describe a hand-scheduled Cooley–Tukey
implementation targeting WH~B0 that achieves 25–75\,GFLOPs/s at
$N \leq 16{,}384$ using multi-core pipelines. Their work focuses on
power-of-two $N$ below the CB overflow boundary. By contrast, our library
extends support to arbitrary $N$ through multi-pass decomposition and
Bluestein's algorithm, at the cost of lower per-transform throughput at small
$N$. The two libraries are complementary: at power-of-two $N \leq 16{,}384$,
Brown et al.'s kernel is faster; at $N > 16{,}384$ or non-power-of-two $N$,
\texttt{ttnn.experimental.fft} is the only available implementation on WH~B0.

\textbf{General FFT algorithm frameworks.}
FFTW~\cite{fftw} pioneered auto-tuned plan-based FFT on CPUs. Spiral~\cite{spiral}
extends this to automatic SIMD code generation. Neither targets accelerators
with per-core private SRAM constraints. Our algorithm-router plays an analogous
role to the FFTW planner but operates at kernel-dispatch granularity with fixed,
pre-compiled kernels rather than generated code.

\textbf{Non-power-of-two transforms.}
Bluestein's chirp-z algorithm~\cite{bluestein} enables arbitrary-$N$ DFT via
two power-of-two FFTs of length $M \geq 2N$. Our contribution is the first
integration of Bluestein into a batched, hardware-aware FFT library for a
neuromorphic-style AI accelerator, including correct \texttt{fp32}/\texttt{bf16}
chirp pre-computation and padding to 1\,024-alignment to satisfy the Tensix
RISC kernel constraints.

\section{Conclusion}
\label{sec:conclusion}

We have presented \texttt{ttnn.experimental.fft}, the first unified
arbitrary-$N$ 1-D FFT library for the Tenstorrent Wormhole~B0 AI accelerator.
The library solves the fundamental circular-buffer overflow barrier inherent to
multi-pass FFTs on private-L1 architectures through \texttt{rebank\_rm}, a new
hardware primitive that reshapes tensors in L1 without a page-sized intermediate
buffer. Four algorithm tiers — Stockham, two-pass, three-pass, and Bluestein —
cover all $N$ up to $16{,}777{,}216$ (power-of-two) and arbitrary non-power-of-two
$N$ for the first time on this platform.

Key results from evaluation on WH~B0:
\begin{itemize}
  \item Peak throughput of 3.67\,GFLOPs/s (Stockham, $N=1{,}024$, $B=64$)
        and 1.84\,GFLOPs/s for single large-$N$ transforms (two-pass,
        $N=1{,}048{,}576$).
  \item Native \texttt{bfloat16} support with only 5–7\% wall-clock
        overhead versus \texttt{fp32}.
  \item Relative L2 accuracy $\leq 4.2{\times}10^{-4}$ for \texttt{fp32}
        across all tiers.
  \item Energy per transform as low as 21\,mJ for the smallest Stockham case
        and scaling predictably to 47\,J for $N = 16{,}777{,}216$.
\end{itemize}

\textbf{Future work.}
Streaming kernels for Bluestein trim/pad would remove the current 1\,024-alignment
constraint and unlock non-power-of-two $N > 65{,}536$. Extending the library to
2-D and batched multi-dimensional FFTs, and to Blackhole and future Tenstorrent
architectures, is a natural next step. Kernel-level double-buffering between DRAM
and L1 could reduce the penalty of the \texttt{rebank\_rm} pass for three-pass
transforms at very large $N$.

% ══════════════════════════════════════════════════════════════════════════════
% REFERENCES
% ══════════════════════════════════════════════════════════════════════════════
\bibliographystyle{IEEEtran}
\bibliography{refs}

\end{document}
